%%%%%%%%%%%%%%%%%%%%%%%%%%%%%%%%%%%%%%%%%%%%%%%%%%%%%%%%%%%%%%%%%%%%%%%%%%%%%%%%
%% 
\cleardoubleoddpage%  Make sure to start each chapter on a new odd page
\chapter{REST-Schnittstellen}
\begin{table}[h]
	\centering
	\caption{REST-Schnittstellen}\label{tab:rest-api}
	\begin{tabularx}{\linewidth}{l >{\raggedright\arraybackslash}X >{\raggedright\arraybackslash}X}
		\toprule
		\textbf{Methode} & \textbf{Pfad} & \textbf{Beschreibung} \\
		\midrule
		\multicolumn{3}{l}{\textit{Schedule}} \\
		\midrule
		\texttt{POST}   & \texttt{/api/v1/schedule}                          & Dienstplan neu erstellen \\
		\texttt{POST}   & \texttt{/api/v1/schedule/reschedule}               & Dienstplan neu berechnen \\
		\midrule
		\multicolumn{3}{l}{\textit{Skills}} \\
		\midrule
		\texttt{GET}    & \texttt{/api/v1/skills}                            & Alle Skills laden \\
		\texttt{POST}   & \texttt{/api/v1/skills}                            & Neuen Skill anlegen \\
		\texttt{DELETE} & \texttt{/api/v1/skills/:id}                        & Skill löschen \\
		\midrule
		\multicolumn{3}{l}{\textit{Station Assignments}} \\
		\midrule
		\texttt{GET}    & \texttt{/api/v1/station-assignments/all}           & Alle Zuweisungen laden \\
		\texttt{GET}    & \texttt{/api/v1/station-assignments/:date}               & Zuweisungen nach Datum filtern \\
		\texttt{PUT}    & \texttt{/api/v1/station-assignments/replace/:date} & Zuweisungen eines Datums ersetzen \\
		\midrule
		\multicolumn{3}{l}{\textit{Stations}} \\
		\midrule
		\texttt{GET}    & \texttt{/api/v1/stations}                          & Alle Stationen laden \\
		\texttt{GET}    & \texttt{/api/v1/stations/:id}                      & Station nach ID laden \\
		\texttt{POST}   & \texttt{/api/v1/stations}                          & Neue Station anlegen \\
		\texttt{PUT}    & \texttt{/api/v1/stations/:id}                      & Station aktualisieren \\
		\texttt{DELETE} & \texttt{/api/v1/stations/:id}                      & Station löschen \\
		\midrule
		\multicolumn{3}{l}{\textit{Users}} \\
		\midrule
		\texttt{GET}    & \texttt{/api/v1/users}                             & Alle Benutzer laden \\
		\texttt{GET}    & \texttt{/api/v1/users/:id}                         & Benutzer nach ID laden \\
		\texttt{POST}   & \texttt{/api/v1/users}                             & Neuen Benutzer anlegen \\
		\texttt{PUT}    & \texttt{/api/v1/users/:id}                         & Benutzer aktualisieren \\
		\texttt{DELETE} & \texttt{/api/v1/users/:id}                         & Benutzer löschen \\
		\bottomrule
	\end{tabularx}
\end{table}

\section{Data Transfer Objects}
Die Kommunikation zwischen Frontend und Backend basiert auf klar definierten Data Transfer Objects (DTOs). 
Im Folgenden werden die wichtigsten Strukturen erläutert.

\begin{lstlisting}[language=Python, caption={User \& UserConstraint}]
	interface User {
		id: number;
		name: string;
		skills: string[];
	}
	
	interface UserConstraint {
		userId: string;
		maxDaysPerMonth?: number | null;
		minDaysPerMonth?: number | null;
		exactDaysPerMonth?: number | null;
		fixedDays?: number[] | null;
		blockedDays?: number[] | null;
	}
\end{lstlisting}

Ein Benutzer besitzt eine Liste von Skills. \texttt{UserConstraint} definiert 
die individuellen Planungsrestriktionen, etwa Mindest- und Maximaleinsatztage 
oder gesperrte Tage.

\begin{lstlisting}[language=Python, caption={Station \& StationAssignment}]
	interface Station {
		id: number;
		name: string;
		maxAssignments: number;
		skills_needed: string[];
	}
	
	interface StationAssignment {
		id: number;
		date: string; // YYYY-MM-DD
		station_id: number;
		user_id: number;
		user: { id: number; name: string };
	}
\end{lstlisting}

Eine Station gibt an, wie viele Personen maximal zugewiesen werden dürfen und 
welche Skills dafür erforderlich sind. \texttt{StationAssignment} verknüpft 
einen Benutzer mit einer Station an einem konkreten Datum.

\begin{lstlisting}[language=Python, caption={SchedulePayload}]
	interface SchedulePayload {
		currentMonth: number;
		currentYear: number;
		daysInMonth: number;
		keepExistingAssignments: boolean;
		newPlan: boolean;
		alternativePlan: boolean;
		constraints: UserConstraint[];
	}
\end{lstlisting}

\texttt{SchedulePayload} wird beim Erstellen oder Neuberechnen eines Dienstplans 
an das Backend übermittelt. Die Flags steuern das Planungsverhalten, die 
\texttt{constraints} bündeln die individuellen Vorgaben aller Benutzer.


%% 
%%%%%%%%%%%%%%%%%%%%%%%%%%%%%%%%%%%%%%%%%%%%%%%%%%%%%%%%%%%%%%%%%%%%%%%%%%%%%%%%
